\chapter*{Introduction générale}
\addcontentsline{toc}{chapter}{Introduction générale}
\markboth{Introduction générale}{}

De nos jours, le domaine du bien-être et de la santé suscite un intérêt grandissant, incitant un nombre croissant de personnes à adopter une activité physique régulière. Cependant, malgré cet engouement, de nombreux pratiquants se heurtent à un problème majeur : le manque de suivi personnalisé. En l'absence d'un encadrement adéquat, les risques de blessures liés à des postures incorrectes augmentent, et la motivation finit souvent par chuter face à l'absence de résultats tangibles ou de programmes inadaptés. Les coachs personnels, bien qu'efficaces, restent une solution onéreuse et peu accessible pour la majorité.

Face à ce constat, une question essentielle se pose : \textbf{Comment exploiter les avancées de l'Intelligence Artificielle pour offrir une expérience de coaching sportif personnalisée, accessible à tous et capable d'accompagner l'utilisateur en temps réel via son smartphone ?}

C'est dans ce contexte que s'inscrit notre projet de fin d'année au sein de l'École Nationale d'Intelligence Artificielle de Berkane (ENIAD). L'objectif principal de notre travail est de concevoir et développer une application mobile de fitness intelligent de nouvelle génération. Pour répondre à la problématique soulevée, ce projet vise concrètement à :
\begin{itemize}
    \item Intégrer des modèles d'apprentissage automatique (Machine Learning) pour analyser les performances et proposer des recommandations d'entraînement et de nutrition sur mesure.
    \item Concevoir une architecture logicielle robuste et évolutive adaptée aux besoins des utilisateurs mobiles.
    \item Offrir une interface utilisateur (UI/UX) ergonomique, intuitive et motivante, afin d'encourager la régularité et l'atteinte des objectifs de santé.
\end{itemize}

Afin de présenter de manière claire et détaillée le fruit de notre travail, ce présent rapport est structuré en six chapitres :
\begin{itemize}
    \item Le \textbf{Chapitre 1} dresse une présentation de notre établissement, l'ENIAD de Berkane.
    \item Le \textbf{Chapitre 2} est consacré à la définition de notre mission, détaillant la problématique, les besoins fonctionnels et les objectifs du projet.
    \item Le \textbf{Chapitre 3} expose les différents outils, langages et frameworks technologiques qui ont été choisis pour mener à bien ce projet.
    \item Le \textbf{Chapitre 4} détaille la phase de réalisation du projet, de la conception architecturale jusqu'à l'implémentation de la solution et des algorithmes d'intelligence artificielle.
    \item Le \textbf{Chapitre 5} dresse l'état actuel du projet et met en lumière les perspectives d'évolution envisageables à court et long terme.
    \item Enfin, le \textbf{Chapitre 6} clôturera cette étude en revenant sur les principales difficultés rencontrées au cours du développement et les solutions adoptées pour les surmonter.
\end{itemize}